% rppg-selective-arrhythmia — IEEE J-BHI draft
%
% Compile:  pdflatex main.tex; bibtex main; pdflatex main.tex; pdflatex main.tex
% Requires: texlive-publishers (for IEEEtran), texlive-bibtex-extra
%
% Class:    IEEEtran journal (single-column draft mode for first reviews;
%           switch to default 2-column for submission by removing the
%           [draftcls,onecolumn] options).

\documentclass[journal]{IEEEtran}

\usepackage[utf8]{inputenc}
\usepackage[T1]{fontenc}
\usepackage{microtype}
\usepackage{amsmath,amssymb,amsfonts}
\usepackage{algorithmic}
\usepackage{graphicx}
\usepackage{textcomp}
\usepackage{xcolor}
\usepackage{booktabs}
\usepackage{multirow}
\usepackage{array}
\usepackage{cite}
\usepackage[hidelinks]{hyperref}
\usepackage{url}

\def\BibTeX{{\rm B\kern-.05em{\sc i\kern-.025em b}\kern-.08em
    T\kern-.1667em\lower.7ex\hbox{E}\kern-.125emX}}

\begin{document}

\title{Learned Class-Conditional Signal-Quality Deferral for Selective rPPG-Based Atrial Fibrillation Screening}

\author{Muhammad~Shahnawaz~Khan%
\thanks{M. S. Khan is an independent researcher, Lahore, Pakistan
        (e-mail: shahnawaz.jrw@gmail.com; ORCID: 0009-0007-4055-6563).}%
\thanks{Manuscript drafted June 2026.}}

\markboth{IEEE Journal of Biomedical and Health Informatics, Submission Draft}%
{Khan: Learned Class-Conditional Signal-Quality Deferral for rPPG-AF Screening}

\maketitle

\begin{abstract}
Contactless atrial fibrillation (AF) screening from facial-video remote photoplethysmography (rPPG) is an attractive opportunistic-screening modality, but deployment requires the model to abstain on segments it cannot classify safely. Standard selective-prediction methods rank samples by model confidence alone, ignoring the physical signal-quality information available from the rPPG pulse waveform itself. We study post-hoc deferral policies that combine model confidence with spectral signal-to-noise ratio (SNR) in the cardiac band. We first show that a na\"{i}ve linear combination achieves an apparent 18\,\% relative improvement in area under the risk-coverage curve (AURC) over uncertainty-only deferral on a 45{,}064-segment synthetic-rPPG benchmark derived from the PhysioNet/CinC 2017 AF Challenge --- but this aggregate gain is clinically harmful: AF recall at 50\,\% coverage collapses from 0.71 to 0.04 because AF's irregular rhythm is itself a low-SNR signature in the cardiac band, so the deferral rule systematically rejects positives. We then introduce \textbf{Learned Class-Conditional Signal-Quality Deferral (LW-CCSD)}, a constrained-grid-search procedure that estimates per-predicted-class quality weights on validation data subject to a configurable AF-recall floor. LW-CCSD makes the safety / accuracy trade-off explicit and tunable: across the Pareto frontier we measure +4.1\,\% AURC improvement at $\leq 0.01$ AF-recall cost, scaling to +15.6\,\% AURC at 8 percentage-point AF cost. The method generalises across three uncertainty quantification methods (deterministic single-pass softmax, MC Dropout, deep ensembles). Strikingly, LW-CCSD applied to a single-forward-pass classifier achieves better selective performance than a five-model deep ensemble without LW-CCSD, suggesting the method can substitute for compute-expensive ensembling in deployed clinical screens.
\end{abstract}

\begin{IEEEkeywords}
Remote photoplethysmography (rPPG), atrial fibrillation, selective prediction, uncertainty quantification, conformal prediction, calibration, contactless cardiac screening, signal quality, deferral policies.
\end{IEEEkeywords}

% =============================================================================
\section{Introduction}
% =============================================================================
\IEEEPARstart{A}{trial} fibrillation (AF) is the most common sustained cardiac arrhythmia and a major risk factor for stroke. It is frequently asymptomatic and paroxysmal, so opportunistic screening outside clinical settings has high public-health value. Remote photoplethysmography (rPPG), which recovers a pulse waveform from skin-colour variation in face video, has emerged as a promising contactless modality. Recent studies report high accuracy for rPPG-based AF detection. However, none report risk-coverage curves, calibration error, or principled deferral policies. rPPG signal quality varies sharply with motion, lighting, skin tone, frame rate, and camera quality, so a confident prediction on a low-quality signal is not deployable.

This work formalises rPPG-based AF screening as a \emph{selective-prediction} problem: a system that defers cases it cannot classify safely. We make three contributions:

\begin{enumerate}
\item We benchmark three uncertainty-quantification (UQ) methods --- MC Dropout, deep ensembles, and conformal prediction --- on a 45{,}064-segment synthetic-rPPG substrate derived from the PhysioNet/CinC 2017 AF Challenge. The best UQ method is dataset-dependent, with consistent rankings across two real ECG datasets.

\item We identify and explain a clinically-harmful failure mode in na\"{i}ve signal-quality-aware deferral: combining model confidence with spectral SNR boosts aggregate AURC by 18\,\% but collapses AF recall from 0.71 to 0.04 at 50\,\% coverage, because AF's irregular rhythm is itself a low-SNR signature.

\item We introduce Learned Class-Conditional Signal-Quality Deferral (LW-CCSD), a post-hoc, model-agnostic procedure that learns per-predicted-class quality weights subject to a configurable AF-recall floor. LW-CCSD makes the safety / accuracy trade-off explicit and tunable, recovering 80\,\% of the na\"{i}ve aggregate gain while preserving AF recall.
\end{enumerate}

A deployment-relevant corollary: LW-CCSD applied to a single-pass deterministic classifier outperforms an unaugmented 5-model deep ensemble. Lightweight inference plus the free post-hoc deferral rule can replace compute-expensive ensembling.

% =============================================================================
\section{Related Work}
% =============================================================================
\subsection{rPPG-Based AF Detection}
% Li et al. 2018 (OBF database), Yan et al. 2022 (Nature Sci Rep), Egorov et al. 2025 (MCD-rPPG).
% Prior rPPG-AF work reports point accuracy without selective-prediction analysis.
[TODO: 200--250 words]

\subsection{Uncertainty Quantification for Medical AI}
% Gal & Ghahramani 2016 (MC Dropout), Lakshminarayanan et al. 2017 (Ensembles),
% Sensoy et al. 2018 (EDL), Liu et al. 2020 (SNGP), Angelopoulos & Bates 2021 (Conformal).
% Calibration: Naeini et al. 2015 (ECE). Selective: Geifman & El-Yaniv 2017.
[TODO: 200--250 words]

\subsection{Signal-Quality-Aware Deferral}
% PPG/rPPG signal-quality literature treats SQI as a hard gate. We treat it as a
% continuous deferral feature combined with UQ.
% Author's prior OACSP work on ordinal-aware class-conditional selective prediction
% in DR grading establishes the class-conditional pattern in a different domain.
[TODO: 150--200 words]

% =============================================================================
\section{Methods}
% =============================================================================
\subsection{Synthetic rPPG from ECG}
% MIT-BIH and CinC 2017 ECG -> R-peak detection (neurokit2) -> asymmetric Gaussian beat
% template at PTT lag -> resample to 30 Hz (camera fps) -> Gaussian noise + baseline wander.
% This is a controlled experiment for the rhythm component of AF detection.
[TODO: methods + equation for beat synthesis]

\subsection{Classifier and Uncertainty Heads}
% 1D-CNN + Transformer over the synthesised pulse. 3-class CE with class weights.
% UQ heads: MC Dropout T=30, deep ensembles M=5, split conformal alpha=0.1.
[TODO: architecture + training details]

\subsection{Selective Evaluation}
% Risk-coverage curves; AURC; selective accuracy at coverage 0.5, 0.7, 0.8, 0.9, 0.95, 1.0;
% ECE (15 bins); multi-class Brier score; per-class recall at coverage.
[TODO: definitions + equations]

\subsection{Signal-Quality Features}
% Spectral SNR (in-band power 0.7--3 Hz vs out-of-band power); template SQI (mean
% per-beat correlation to median template). Both computed deterministically from
% the 30 Hz pulse waveform at inference time.
[TODO: equations]

\subsection{Class-Conditional Deferral Policies}
We compare four scoring functions for selective prediction. Let $\hat{p}(x)$ denote model confidence, $q(x)$ the signal-quality feature (rank-normalised within the test set), and $\hat{c}(x)$ the predicted class.

\begin{align}
\text{UQ-only:}\quad & s_0(x) = \text{rank}(\hat{p}(x)) \\
\text{Na\"{i}ve SQI:}\quad & s_w(x) = (1-w)\,\text{rank}(\hat{p}(x)) + w\,\text{rank}(q(x)) \\
\text{AF-immune:}\quad & s_w^{\text{AI}}(x) =
  \begin{cases}
    \text{rank}(\hat{p}(x)) & \hat{c}(x) = \text{AF} \\
    s_w(x) & \text{otherwise}
  \end{cases} \\
\text{LW-CCSD:}\quad & s_{\mathbf{w}}^{\text{LW}}(x) = (1-w_{\hat{c}(x)})\,\text{rank}(\hat{p}(x)) + w_{\hat{c}(x)}\,\text{rank}(q(x))
\end{align}

In LW-CCSD, $\mathbf{w} = (w_{\text{NSR}}, w_{\text{AF}}, w_{\text{Other}})$ is learned by constrained grid search on the validation set:
\begin{align}
\mathbf{w}^{*} = \arg\min_{\mathbf{w} \in [0,1]^3}\ &\text{AURC}_{\text{val}}(s_{\mathbf{w}}^{\text{LW}}) \\
\text{s.t.}\quad & \text{recall}_{\text{val}}^{(c)}(s_{\mathbf{w}}^{\text{LW}}; c_0) \geq f_c\ \forall c
\end{align}
where $\text{recall}_{\text{val}}^{(c)}(s; c_0)$ denotes class-$c$ recall on val at coverage $c_0$, and $f_c$ is the per-class recall floor.

% =============================================================================
\section{Experimental Setup}
% =============================================================================
\subsection{Datasets}
% MIT-BIH Arrhythmia DB (47 records, 3-class NSR/AF/Other) for pipeline validation.
% PhysioNet/CinC 2017 AF Challenge (8,528 single-lead ECGs, 738 AF) for the main
% benchmark. Records split 70/15/15 stratified by class, subject-disjoint.
[TODO: counts, splits, sizes]

\subsection{Implementation}
% PyTorch 2.x; Apple Silicon MPS; CUDA-portable.
% Open-source release: https://github.com/ShahnawazKakarh/rppg-selective-arrhythmia
[TODO]

\subsection{Evaluation Protocol}
% Validation set used for conformal calibration AND for LW-CCSD weight selection.
% No test-set tuning. AF-recall floor at constraint coverage = 0.5 (most adversarial).
[TODO]

% =============================================================================
\section{Results}
% =============================================================================
\subsection{UQ Benchmark on Synthetic rPPG}
% Three UQ methods on 45,064-segment CinC synth set.
% MC Dropout AURC 0.198, Conformal 0.237, Ensembles 0.215. Best method depends on metric.

\begin{table}[t]
\centering
\caption{UQ benchmark on synth-rPPG (CinC test, $n{=}6{,}750$). Best per column in \textbf{bold}.}
\label{tab:uq_comparison}
\begin{tabular}{lrrrr}
\toprule
Method & Acc & ECE $\downarrow$ & Brier $\downarrow$ & AURC $\downarrow$ \\
\midrule
MC Dropout (T{=}30) & \textbf{0.715} & 0.076 & 0.418 & \textbf{0.198} \\
Conformal ($\alpha{=}0.1$) & 0.699 & \textbf{0.056} & 0.439 & 0.237 \\
Ensembles (M{=}5) & 0.686 & 0.064 & 0.440 & 0.215 \\
\bottomrule
\end{tabular}
\end{table}

\subsection{Na\"{i}ve SQI Deferral: An 18\,\% Gain That Is Clinically Harmful}
% Aggregate AURC 0.197 vs 0.237 UQ-only baseline (+18%), but AF recall collapses
% from 0.707 to 0.043 at 50% coverage. Mechanism: AF rhythm = low SNR.

\begin{table}[t]
\centering
\caption{Na\"{i}ve SQI deferral. Aggregate AURC improves but AF recall collapses.}
\label{tab:naive_sqi}
\begin{tabular}{lrrrr}
\toprule
Coverage & UQ AURC & SQI AURC & UQ AF recall & SQI AF recall \\
\midrule
0.50 & --- & --- & 0.707 & 0.043 \\
0.70 & --- & --- & 0.732 & 0.365 \\
0.80 & --- & --- & 0.764 & 0.637 \\
0.95 & --- & --- & 0.814 & 0.778 \\
\bottomrule
\end{tabular}
\end{table}

\subsection{LW-CCSD: The Pareto Frontier}
% Headline figure: floor 0.60 -> +4.1% AURC at -0.004 AF recall.
%                 floor 0.55 -> +15.6% AURC at -0.081 AF recall.

\begin{table}[t]
\centering
\caption{LW-CCSD Pareto frontier on synth-rPPG CinC test set (deterministic classifier).}
\label{tab:pareto}
\begin{tabular}{lcrrr}
\toprule
Val AF floor & $\mathbf{w}^{*}$ (NSR/AF/Other) & Test AURC & $\Delta$ vs UQ & AF cost \\
\midrule
0.60 & 0.0 / 0.1 / 0.5 & 0.2270 & +4.1\,\% & $-$0.004 \\
0.58 & 0.0 / 0.4 / 0.6 & 0.2082 & +12.1\,\% & $-$0.045 \\
0.55 & 0.3 / 0.4 / 0.4 & 0.1998 & +15.6\,\% & $-$0.081 \\
0.50 & 0.1 / 0.5 / 0.6 & 0.2020 & +14.7\,\% & $-$0.123 \\
0.45 & 0.2 / 0.5 / 0.5 & 0.1984 & +16.2\,\% & $-$0.166 \\
0.40 & 0.4 / 0.5 / 0.5 & 0.1958 & +17.3\,\% & $-$0.220 \\
\midrule
UQ-only (baseline) & --- & 0.2367 & --- & ~0.000 \\
Na\"{i}ve SQI (w{=}0.7) & shared & 0.1942 & +18.0\,\% & $-$0.664 \\
\bottomrule
\end{tabular}
\end{table}

\subsection{Cross-UQ Replication}
% LW-CCSD applied to deterministic, MC Dropout, and ensemble confidences.
% Pattern: SQI value scales inversely with UQ informativeness.
% Deployment punchline: LW-CCSD on deterministic > ensemble UQ-only.

\subsection{Mechanism: Why Na\"{i}ve SQI Fails on AF}
% Quantitative per-class SNR distribution analysis.

% =============================================================================
\section{Discussion}
% =============================================================================
% - Why the method works: SQI provides orthogonal information when UQ is weak;
%   the class-conditional weight automatically discovers where it does not.
% - Deployment implications: LW-CCSD replaces ensembles.
% - Limitations: validated on synthetic rPPG; OBF/MAHNOB-HCI replication pending.
% - Generality: the pattern (class-defining signal collides with auxiliary feature)
%   recurs across domains (cf. the author's OACSP work in DR grading); LW-CCSD
%   is a general fix.
[TODO: 1--1.5 pages]

% =============================================================================
\section{Conclusion}
% =============================================================================
We introduced LW-CCSD, a post-hoc class-conditional deferral procedure that learns per-predicted-class signal-quality weights subject to a configurable safety floor. On a 45{,}064-segment synthetic-rPPG benchmark, LW-CCSD recovers most of the aggregate AURC gain of na\"{i}ve signal-quality deferral while preserving AF recall, and applied to a lightweight single-pass classifier it outperforms an unaugmented five-model deep ensemble. The method is model-agnostic, requires no retraining, and generalises across three uncertainty quantification methods. We release the full pipeline and CinC-derived synth-rPPG substrate under MIT licence; replication on real face-video AF data is a natural next step pending OBF / MAHNOB-HCI access.

% =============================================================================
\section*{Acknowledgments}
% =============================================================================
This work uses public datasets --- PhysioNet/CinC 2017 AF Challenge and the MIT-BIH Arrhythmia Database --- whose maintainers are gratefully acknowledged.

% =============================================================================
\section*{Data and Code Availability}
% =============================================================================
Open-source pipeline including all scripts, configurations, trained checkpoints, and per-method evaluation artefacts is available at \url{https://github.com/ShahnawazKakarh/rppg-selective-arrhythmia} under MIT licence. The CinC-derived synthetic-rPPG cache is reproducible from the released code in approximately 20 minutes of CPU time on a recent laptop.

% =============================================================================
\section*{Conflict of Interest Statement}
% =============================================================================
The author declares no competing interests.

% =============================================================================
\bibliographystyle{IEEEtran}
\bibliography{refs}

\end{document}
